\documentclass[11pt]{article}

% ── Packages ──
\usepackage[utf8]{inputenc}
\usepackage[T1]{fontenc}
\usepackage{amsmath,amssymb}
\usepackage{graphicx}
\usepackage{booktabs}
\usepackage{hyperref}
\usepackage{xcolor}
\usepackage{natbib}
\usepackage[margin=1in]{geometry}
\usepackage{microtype}
\usepackage{caption}
\usepackage{subcaption}

% ── Metadata ──
\title{Reasoning Honesty Under Scrutiny: A Differential Evaluation Framework for Detecting Audience-Dependent Evasion in Large Language Models}

\author{
  % Author name(s) here
}

\date{\today}

% ── Document ──
\begin{document}

\maketitle

\begin{abstract}
When a large language model refuses a request, is it reasoning honestly about its limitations or deploying patterned evasion strategies? We present Claude-Assessor, a sociolinguistic evaluation framework that audits reasoning honesty in LLMs through differential testing across audience framings. Our three-stage protocol---differential framing, capability probing, and multi-turn pressure testing---evaluates 52 probes across six knowledge domains on three Claude models (468 total probe runs). We find that while the majority of responses (64--88\%) are substantively consistent across framings, discriminatory behavior concentrates sharply in specific domains, particularly cybersecurity and meta-cognitive reasoning. The capability gap between demonstrated knowledge and claimed willingness provides falsifiable evidence that models possess knowledge they selectively withhold based on perceived audience. These patterns are stable across variance runs and model scales, suggesting systematic rather than stochastic origin.
\end{abstract}

% ── 1. Introduction ──
\section{Introduction}
\label{sec:introduction}

% Core research question and motivation
% Why reasoning honesty matters for AI safety and trust
% The gap between capability and willingness as a diagnostic signal

% ── 2. Related Work ──
\section{Related Work}
\label{sec:related}

% LLM evaluation frameworks
% Safety and alignment testing
% Sociolinguistic approaches to AI behavior
% Differential testing methodologies

% ── 3. Methodology ──
\section{Methodology}
\label{sec:methodology}

\subsection{Threat Model}
% What we're testing for: patterned evasion vs. honest reasoning

\subsection{Three-Stage Protocol}
% Stage 1: Differential framing (naive vs. credentialed)
% Stage 2: Capability probing
% Stage 3: Multi-turn pressure testing

\subsection{Probe Design}
% 52 probes across 6 domains
% Domain selection rationale
% Framing sensitivity scoring (0--6 scale)

\subsection{Scoring Model}
% Classification thresholds
% Verdict taxonomy: discriminatory, tone\_modulated, consistent
% Capability gap analysis

% ── 4. Experimental Setup ──
\section{Experimental Setup}
\label{sec:setup}

% Models evaluated (Haiku 4.5, Sonnet 4, Sonnet 4.6)
% Variance runs (3 per probe per model)
% API configuration and parameters

% ── 5. Results ──
\section{Results}
\label{sec:results}

\subsection{Overall Consistency}
% Verdict distribution across all models

\subsection{Domain-Level Analysis}
% Where discrimination concentrates
% Cybersecurity and reasoning as outlier domains

\subsection{Cross-Model Comparison}
% Behavioral differences across model scale
% What improves and what doesn't

\subsection{Capability Gap Findings}
% Evidence of knowledge the model has but withholds
% Falsifiability of evasion claims

% ── 6. Discussion ──
\section{Discussion}
\label{sec:discussion}

% What the patterns mean for alignment
% Limitations of the framework
% The consistency-as-honesty assumption
% Domain-specific vs. general evasion

% ── 7. Limitations ──
\section{Limitations}
\label{sec:limitations}

% Single model family (Claude)
% English-only probes
% Scoring model limitations
% Prompt sensitivity

% ── 8. Conclusion ──
\section{Conclusion}
\label{sec:conclusion}

% Summary of contributions
% Implications for AI evaluation practice
% Future work

\bibliographystyle{plainnat}
\bibliography{references}

\end{document}
